\begin{abstract}
Knowledge-graph quality control couples two decisions: which constraints to acquire and which graph defects to repair under a limited budget. We formulate this interaction as a finite-horizon control problem and train a Double-DQN policy over four graph operations and four rule-management operations. A dual-strategy module proposes text-derived rule candidates through deletion completion and augmentation expansion. The reward measures introduced violations by identity after each edit. On a controlled-corruption benchmark, ten retrained Double-DQN policies reach mean normalized final quality 0.96855 and introduce no new violations. An acquire-then-deficit heuristic reaches 0.98248 but introduces 4.8 new violations per run; mean discounted returns are 0.3660 and 0.3508, respectively. Forty retrained ablation models support the action mask's contribution to final quality, while other component gains do not pass multiplicity correction. Replaying 9,456 archived responses shows complementary candidates on the same documents, but augmentation alone yields more constraint candidates at equal call budgets. An execution audit compiles 18,143 exact typed patterns; their union detects 10 of 64 defects in the designed RuleTest-94 suite. A separate offline bridge activates these rules and removes ten prohibited records while retaining all thirty clean records. The results characterize the tradeoff between structural improvement and new violations, and distinguish candidate diversity, executable coverage and factual restoration.
\end{abstract}
